% Einheit 1 von 3 – Punkte spiegeln (bank/spiegelung/e1.jsonl, zusammenbau v0.1)
%% TODO Zeitmarke Einheit 1: Marken-Zeile nicht lesbar
%% TODO Zweigzeile Teil 1: was hier gelernt wird (Satz in Schülersprache aus dem Einheitstitel) – Einheit 1
\einheitenkopf[e1]{Einheit 1 von 3 · Punkte spiegeln}
\zweigzeile{Abitur GK}
\setcounter{aufgabe}{8}

%% TODO Titel: Ich-kann-Satz fehlt in der Bank (Platzhalter: Kettenname) – Nr. 9
%% TODO Anweisung: ein Satz über der ersten Teilaufgabe fehlt in der Bank – Nr. 9
\begin{aufgabe}{Was ist gegeben, was gesucht?}
\begin{teile}
\teil Was ist gesucht? „$P$ wird an der Ebene $E$ gespiegelt. Bestimme den Spiegelpunkt.“ \\ \kreuz{das Bild gesucht – spiegeln} \\ \kreuz{die Spiegelfläche gesucht – rückwärts}
\end{teile}
\end{aufgabe}

%% TODO Titel: Ich-kann-Satz fehlt in der Bank (Platzhalter: Kettenname) – Nr. 10
%% TODO Anweisung: ein Satz über der ersten Teilaufgabe fehlt in der Bank – Nr. 10
\begin{aufgabe}{Welche Koordinate wechselt?}
\begin{teile}
\teil Welche Koordinate wechselt? $P(2 | 3 | 5)$ und $P'(2 | 3 | -5)$ \\ \kreuz{symmetrisch zur $x_1x_2$-Ebene} \\ \kreuz{symmetrisch zur $x_1x_3$-Ebene} \\ \kreuz{symmetrisch zur $x_2x_3$-Ebene} \\ \kreuz{zu keiner Koordinatenebene}
\end{teile}
\end{aufgabe}

%% TODO Titel: Ich-kann-Satz fehlt in der Bank (Platzhalter: Kettenname) – Nr. 11
%% TODO Anweisung: ein Satz über der ersten Teilaufgabe fehlt in der Bank – Nr. 11
\begin{aufgabe}{Punkte spiegeln}
%% TODO \rechenplatz{n}: Schrittzahl des Grundfalls fehlt in der Bank (nur bei fester Schreibform, 2.3 b)
\begin{teile}
\teil Was wird gespiegelt – und woran? „Spiegle $A$ am Mittelpunkt $M$ der Strecke $BC$.“ \\ \kreuz{Punkt an Punkt} \\ \kreuz{Punkt an Ebene} \\ \kreuz{Gerade an Ebene} \\ \kreuz{Gerade auf Gerade} \\ \kreuz{Körper an Symmetrieebene}
\teil Spiegle $P(1 | 3 | 0)$ am Punkt $Q(2 | 4 | 2)$. $P'($ \leerfeld | \leerfeld | \leerfeld $)$
\teil Spiegle $P(2 | -3 | 4)$ an der $x_1x_2$-Ebene. $P'($ \leerfeld | \leerfeld | \leerfeld $)$
\teil $F(1 | 2 | 2)$ ist der Lotfußpunkt von $Q(1 | 2 | 5)$ auf $E: x_3 = 2$. Spiegle $Q$ an $E$. $R($ \leerfeld | \leerfeld | \leerfeld $)$
\teil Spiegle $P(4 | 4 | 4)$ an der Ebene $E: x_1 + x_2 + x_3 = 3$. $P'($ \leerfeld | \leerfeld | \leerfeld $)$
\teil $P(1 | 0 | 2)$ wird an einer Ebene $E$ auf $Q(5 | 4 | 0)$ gespiegelt. $R$ und $S$ liegen auf der Geraden $PQ$ symmetrisch zu $E$, $R$ auf der Seite von $P$, und es gilt $|RS| = 2 \cdot |PQ|$. Bestimme $R$. \hfill (Abitur 2022 LK) \\ $R($ \leerfeld | \leerfeld | \leerfeld $)$
\end{teile}
\end{aufgabe}

%% TODO Titel: Ich-kann-Satz fehlt in der Bank (Platzhalter: Kettenname) – Nr. 12
%% TODO Anweisung: ein Satz über der ersten Teilaufgabe fehlt in der Bank – Nr. 12
\begin{aufgabe}{Punkte spiegeln – Fehler finden, Begründen}
\begin{teile}
\teil Finde den Fehler. Spiegle $P(2 | 1 | 3)$ am Punkt $Z(3 | 4 | 1)$. \rechnung{\overrightarrow{PZ} &= (1 | 3 | -2) \\ P' &= P + \overrightarrow{PZ} = (3 | 4 | 1)}
\teil Warum wechselt beim Spiegeln an einer Koordinatenebene genau ein Vorzeichen?
\end{teile}
\end{aufgabe}
